\documentclass[11pt]{article}

\usepackage[margin=25mm]{geometry}

\usepackage{amsmath,amssymb,fontspec,unicode-math}

\IfFontExistsTF{DejaVu Serif}{\setmainfont{DejaVu Serif}}{\setmainfont{Latin Modern Roman}}

\IfFontExistsTF{Latin Modern Math}{\setmathfont{Latin Modern Math}}{\setmathfont{DejaVu Math TeX Gyre}}

\usepackage{xurl,hyperref}

\setlength{\parindent}{0pt}\setlength{\parskip}{6pt}

\begin{document}

\title{Composing Linear Layers from Irreducibles}

\author{ELVIA preliminary AI analysis}\date{21 October 2026}\maketitle

Preliminary assessment: Minor.\par

Independent scientific review remains pending. Author adoption has not been requested.

The full Clifford-action determinant loses orientation information, and two release differentiation/boundary paths need correction. The implemented even-rotor forward family is not shown to contain the forbidden reflection.

\section*{Review scope}

Latest-version, claim-focused review of the cited author-paper passages and applicable current author-release paths. The stated counterexamples and scope controls were checked analytically; no author program, old verifier, training run or benchmark was executed.

\section*{Problem 1: Full-action orientation does not identify the spin image}

Source: Appendix B, Property 3, pp. 27–28.

Assessment: supported; impact: minor

\subsection*{Source claim (paraphrase)}

Property 3 characterizes the spin image using grade preservation and positive determinant on the entire Clifford algebra.

\subsection*{Analysis}

In \(\mathrm{Cl}(2)\), conjugation by \(g=e_1\) has matrix \(\operatorname{diag}(1,1,-1,-1)\) in the scalar/vector/bivector basis. It preserves grades and has full determinant \(+1\), but its grade-one block has determinant \(-1\), a reflection unavailable to a \(\mathrm{Spin}(2)\) rotor. More generally the full exterior lift has determinant \((\det R)^{2^{n-1}}\), losing the sign for \(n\geq2\).

\subsection*{Finding and scope}

The characterization is too broad. The actual even-rotor forward parameterization is not shown to realize that reflection.

\subsection*{Proposed correction}

Impose \(\det(N_g|_{\mathrm{Cl}^1})=1\) on the grade-one action, together with the stated orthogonality and grade conditions.

An explicit local correction is described; author adoption and a corresponding published erratum are not established.

\section*{Problem 2: Detached decomposition omits fixed-point dependence}

Source: Section 4.1, PDF pp.5–6; Appendix B, Theorems 2–3; rotor\_layer.py, Rotor.\_update\_rotors, lines 164–180

Assessment: supported; impact: minor

\subsection*{Source claim (paraphrase)}

The release solves a decomposition without gradient tracking, then reruns with detached warm-start vectors.

\subsection*{Analysis}

For bivector \(b=1.2e_{12}+0.6e_{34}+he_{13}\), the selected matrix entry’s recomputed exponential derivative is approximately \(-0.23271515\), whereas the released backward path gives \(-0.22552789\). Recomputed forward finite differences support the former. Detaching spectral/fixed-point state omits its parameter dependence.

\subsection*{Finding and scope}

The backward rule is not the derivative of the recomputed rotor map on this control. Surrogate-gradient training may still be useful.

\subsection*{Proposed correction}

Track the relevant decomposition iterations or supply a valid implicit derivative under stated nondegeneracy conditions. A local differentiable-iteration control does not prove stability at all spectral collisions.

An explicit local correction is described; author adoption and a corresponding published erratum are not established.

\section*{Problem 3: Zero simple components need an exponential limit and stable shape}

Source: Eq.(5), PDF p.5; current author README prescribed torch\_ga\_fix submission branch at 376a48bfeeba2e76ef42abee6fe516a349d8d3bb, torch\_ga/torch\_ga.py, exp\_simple\_batched

Assessment: supported; impact: minor

\subsection*{Source claim (paraphrase)}

The author README prescribes the torch\_ga\_fix submission branch. Its current exp\_simple\_batched divides a component by its norm before multiplying by sine and squeezes a singleton dimension.

\subsection*{Analysis}

A zero component therefore gives \(0/0\) instead of the identity exponential. Even nonzero simple input \(e_{12}\) in \(\mathrm{Cl}(4)\) leaves a zero second component after decomposition, so zero input is not the only boundary. A singleton batch can additionally lose the axis needed by concatenation.

The dependency’s separate default-main branch contains another API with an explicit zero-angle branch. That is opposing evidence against a claim covering every dependency version; it does not remove the defect from the current author-prescribed submission branch.

\subsection*{Finding and scope}

The helper is undefined on legitimate simple-component inputs. Random dense-training incidence is unknown.

\subsection*{Proposed correction}

Use \(\exp(b)=\cos q+\operatorname{sinc}(q/\pi)b\), \(q=\|b\|\), with normalized sinc value \(1\) at zero, and preserve the batch axis.

The currently prescribed submission branch retains the defect. A different default-main implementation has zero-angle handling, but its integration into the paper release was not established.

\section*{Limits}

Preliminary AI assessment; independent human scientific adjudication remains pending.\par

An analytic counterexample or a release discrepancy does not by itself remeasure the paper’s reported experiments.\par

\section*{Source identity}

arXiv:2507.11688v4 (current author version; updated 2026-06-09T18:47:26Z)

\url{https://arxiv.org/pdf/2507.11688v4}

Reviewed PDF SHA-256: \nolinkurl{a99d414232746689b168daa903ea4046fc535e9c0e55bc237d1b4275f5670ef1}

The selected author/publisher source is the same work as the original case, matched by exact title, author identity and official version channel. The exact selected PDF SHA-256 was recomputed in the latest-version pass for Issue \#27. The PDF is a retained author/publisher copy; its original retrieval time is stated only when supported by matching acquisition metadata. File modification time was not substituted.

\end{document}
